\documentclass[11pt,a4paper]{article}

\usepackage[utf8]{inputenc}
\usepackage[T1]{fontenc}
\usepackage[spanish, es-nodecimaldot]{babel}
\addto\captionsspanish{\renewcommand{\tablename}{Tabla}}
\usepackage{geometry}
\geometry{margin=2.5cm}
\usepackage{booktabs}
\usepackage{array}
\usepackage{hyperref}
\hypersetup{colorlinks=true, linkcolor=black, urlcolor=blue}
\usepackage{parskip}

\title{Informe de mantenimiento — \texttt{org.openmarkov.core}}
\author{Manuel Arias Calleja}
\date{Semana del 16 al 21 de marzo de 2026}

\begin{document}

\maketitle

\section*{Resumen ejecutivo}

Durante esta semana se llevó a cabo una campaña sistemática de mantenimiento y modernización del
módulo \texttt{core}. El trabajo se agrupó en cuatro áreas: modernización de colecciones,
migración del sistema de \emph{logging}, limpieza de código estático y ampliación de la suite de
tests.

\begin{table}[!htbp]
\centering
\begin{tabular}{lcc}
\toprule
\textbf{Área} & \textbf{Commits} & \textbf{Ficheros afectados} \\
\midrule
Colecciones modernas & 2 & 11 \\
\emph{Logging} con log4j2 & 2 & 6 \\
Limpieza de código & 2 & 21 \\
Nuevos tests & 1 & 4 nuevas clases \\
Otros & 1 & 1 \\
\bottomrule
\end{tabular}
\caption{Resumen de actividad por área.}
\end{table}

% ---------------------------------------------------------------
\section{Modernización de colecciones}

\subsection{\texttt{Hashtable} → \texttt{HashMap} / \texttt{Map}}

\emph{20 mar, \texttt{8a44d8c0}.}
Se sustituyeron todos los usos de la clase obsoleta \texttt{java.util.Hashtable} por
\texttt{HashMap} o por la interfaz \texttt{Map} en cinco clases: \texttt{Strategy},
\texttt{StrategyUtilities}, \texttt{SystematicSampling}, \texttt{StrategyTree} y
\texttt{DiscretePotentialOperations}. \texttt{Hashtable} es una clase legada, sincronizada de
forma incondicional y sin soporte para \texttt{null} como clave; \texttt{HashMap} es su sustituto
moderno para contextos de hilo único, que es el caso de todos estos usos.

\subsection{\texttt{Stack} → \texttt{ArrayDeque} / \texttt{Deque}}

\emph{20 mar, \texttt{f9a4c833}.}
Se sustituyeron todos los usos de \texttt{java.util.Stack} por \texttt{ArrayDeque} en seis clases:
\texttt{BasicOperations}, \texttt{ProbNetOperations}, \texttt{Graph}, \texttt{UniqueStack},
\texttt{TreeADDPotential} y \texttt{Util}. \texttt{Stack} extiende \texttt{Vector} (también legado
y sincronizado) e impone una semántica de acceso aleatorio innecesaria para las pilas usadas aquí.

El cambio más delicado fue en \texttt{ProbNetOperations.getSequenceOfDecisions()}, cuyo tipo de
retorno pasó de \texttt{Stack<Variable>} a \texttt{Deque<Variable>}. Para preservar el orden de
iteración compatible con el comportamiento previo se usó \texttt{addLast()} en lugar de
\texttt{push()} (que en \texttt{ArrayDeque} añade por el frente). Se actualizó
\texttt{DANOperations} (módulo \texttt{inference}) para usar \texttt{peekLast()} en lugar del
método \texttt{Stack}-específico \texttt{lastElement()}.

% ---------------------------------------------------------------
\section{Migración del sistema de \emph{logging}}

\emph{21 mar, \texttt{ec670741} + \texttt{4d1f59c3}.}
Se eliminaron todos los usos de \texttt{System.out.println} que quedaban en el módulo y se
reemplazaron por llamadas al \emph{logger} log4j2 con el nivel semántico correcto en seis ficheros:

\begin{table}[!htbp]
\centering
\begin{tabular}{lll}
\toprule
\textbf{Fichero} & \textbf{Situación anterior} & \textbf{Nivel nuevo} \\
\midrule
\texttt{BasicOperations}      & Advertencia: red con >1 nodo de decisión   & \texttt{WARN}  \\
\texttt{TemporalNetOperations}& Copia de enlaces de decisión               & \texttt{DEBUG} \\
\texttt{ExtensionTree}        & Impresión del árbol de extensión           & \texttt{INFO}  \\
\texttt{ProbNet}              & Mensajes de diagnóstico de red             & \texttt{DEBUG} \\
\texttt{Variable}             & Traza de estados de variable               & \texttt{DEBUG} \\
\bottomrule
\end{tabular}
\caption{Sustituciones de \texttt{System.out.println} por log4j2.}
\end{table}

% ---------------------------------------------------------------
\section{Limpieza de código estático}

\subsection{Corrección del orden de modificadores}

\emph{21 mar, \texttt{b6b65634}.}
Se corrigió el orden de los modificadores \texttt{final static} → \texttt{static final} en diez
constantes distribuidas en cinco clases: \texttt{DiscretePotentialOperations} (3 campos),
\texttt{Variable}, \texttt{Node}, \texttt{NetworkTypeUtils} y \texttt{Criterion} (2 campos).
El orden correcto según la especificación del lenguaje Java (JLS §8.3.1) es \texttt{static} antes
de \texttt{final}.

\subsection{Eliminación de código comentado}

\emph{21 mar, \texttt{356fe856}.}
Se eliminaron 290 líneas de código comentado en 16 ficheros. El mayor bloque correspondía a
\texttt{DiscretePotentialOperations} (203 líneas), que contenía implementaciones alternativas
desechadas de operaciones sobre potenciales. El resto se distribuyó entre ficheros de potenciales,
inferencia y localización.

% ---------------------------------------------------------------
\section{Nuevos tests de regresión}

\emph{20 mar, \texttt{9c1f1115}.}
Se añadieron cuatro nuevas clases de test que suman 1\,596 líneas:

\begin{table}[!htbp]
\centering
\begin{tabular}{ll}
\toprule
\textbf{Clase} & \textbf{Cobertura} \\
\midrule
\texttt{ProbNetCopyTest}
  & Copia profunda de redes: nodos, arcos, potenciales, evidencias \\
\texttt{VariableModifyStateTest}
  & Modificación de estados: añadir, renombrar, eliminar, reordenar \\
\texttt{TablePotentialRegressionTest}
  & Operaciones sobre \texttt{TablePotential}: suma, producto, marginalización \\
\texttt{DiscretePotentialOperationsRegressionTest}
  & Operaciones discretas: multiplicación, eliminación de variables \\
\bottomrule
\end{tabular}
\caption{Nuevas clases de test añadidas.}
\end{table}

% ---------------------------------------------------------------
\section{Otras contribuciones}

\textbf{Refactorización de \texttt{ProbNet}} \emph{(18 mar, \texttt{348f87a2}, Jose Andres):}
simplificación menor de dos líneas en la gestión de ediciones.

% ---------------------------------------------------------------
\section*{Estadísticas globales}

\begin{table}[!htbp]
\centering
\begin{tabular}{lr}
\toprule
\textbf{Métrica} & \textbf{Valor} \\
\midrule
Commits                             & 8 \\
Días activos                        & 4 (18–21 mar) \\
Ficheros modificados                & $\approx$30 \\
Líneas eliminadas                   & $\approx$330 \\
Líneas añadidas                     & $\approx$1\,650 \\
Clases de test nuevas               & 4 \\
\bottomrule
\end{tabular}
\caption{Estadísticas globales de la semana.}
\end{table}

\end{document}
